\documentclass{article}
\usepackage{graphicx} % Required for inserting images
\usepackage{amsmath}
\usepackage{amsthm}
\newtheorem{remark}{Remark}
\usepackage{siunitx}
\usepackage[margin=1in]{geometry}
\title{Computations and Assumptions Behind the R-32 Computations in DVCT}
\author{Shilpa Narasimhan}
\date{August 2026}
\begin{document}
\maketitle
\section{Vanilla Python Equations of State Package}\label{sec:vanilla_pkg}
\subsection{PR equation of state for $P,V,T$}\label{sec:pr_pvt}
PR and SRK are both special cases of the generalized cubic equation of state, differing only in their $(u, w)$ parameters and $\alpha$-function coefficients; PR was developed specifically to improve upon SRK's known underprediction of liquid density while retaining comparable vapor-phase accuracy \cite{peng1976,soave1972}. The Peng-Robinson parameterized equation of state, written in the generalized-cubic form adopted by IDAES~\cite{idaes_cubic_eos_2026}, is used:
\begin{equation}\label{eq:PR}
        P = \frac{RT}{V-b} - \frac{a}{V^2 + ubV + wb^2}
\end{equation}
This has an equivalent form in terms of the dimensionless compressibility factor ($Z = \frac{PV}{RT}$):
\begin{subequations}
    \begin{align}
        0&= Z^3 - (1+B-uB)Z^2 + (A-uB-(u-w)B^2)Z - AB - wB^2 - wB^3 \label{eq:cubic_eos}\\
        A& = \frac{a(T) P}{R^2 T^2}\\
        B & = \frac{b P}{RT}
    \end{align}
\end{subequations}
The parameters $a$ and $b$ can be computed as:
\begin{equation}
\begin{aligned}
b &= \frac{\Omega_B RT_c}{P_c}\\
a(T) & = \frac{\Omega_A R^2 T_c^2 \alpha(T)}{P_c}\\
\alpha(T) & = \Bigg[1 + \kappa (1-\sqrt{\frac{T}{T_c}})\Bigg]^2\\
\kappa &= 0.37464 + 1.54226\omega -0.26992\omega^2
\end{aligned}
\end{equation}
The expression for $\kappa$ is obtained from~\cite{peng1976} and here $\omega$ is the acentric factor. The constants $u$, $w$, $\Omega_A$, and $\Omega_B$ are universal (the same for every fluid) and are also obtained from~\cite{peng1976}, written here in the generalized form used by~\cite{idaes_cubic_eos_2026}:
\[
u = 2, \qquad w = -1, \qquad \Omega_A = 0.45724, \qquad \Omega_B = 0.0778
\]
$T_c$, $P_c$, and the acentric factor $\omega$ are fluid-specific and are not universal constants.
Because of the $1/(V-b)$ singularity in Eq.~\eqref{eq:PR}, only real roots satisfying $Z>B$ (equivalently $V>b$, since $b$ represents the finite excluded molecular volume) are physically admissible; complex roots and any real root with $Z \le B$ are discarded. Solving the cubic equation in Eq.~\eqref{eq:cubic_eos} numerically yields either one real root, corresponding to a single-phase state (vapor, liquid, or supercritical), or three real roots, corresponding to a state within the two-phase region. In the former case, the single real root is the unique physically valid compressibility factor at that $(T,P)$, so no root-selection ambiguity exists. However, labeling that state as liquid, vapor, or supercritical is not determined by the value of the root itself, but by comparing $(T,P)$ against the saturation curve: if $T>T_c$, the state is supercritical and no liquid/vapor distinction applies; otherwise, $P>P_{sat}(T)$ corresponds to liquid and $P<P_{sat}(T)$ corresponds to vapor. In the latter case, the smallest and largest real roots, denoted $Z_L$ and $Z_V$ respectively, correspond to the coexisting saturated liquid and vapor, while the middle root is discarded as mechanically unstable ($\frac{\partial P}{\partial V}>0$).
\subsubsection{Acentric factor}\label{sec:acentric_factor}
The acentric factor $\omega$ follows the Pitzer definition~\cite{reid1987}, based on the reduced saturation pressure at a reduced temperature of $T_r=0.7$:
\begin{equation}
    \omega \equiv -\log_{10}\left(P_r^{sat}\right)\Big|_{T_r=0.7} - 1, \qquad P_r^{sat} = \frac{P^{sat}(T_r=0.7 \cdot T_c)}{P_c}
\end{equation}
$P^{sat}$ at $T_r=0.7$ is obtained from the Linde datasheet for R-32, which is shared across all three parameter sets (NIST, GCGP, SPGP) considered in this work. Since $T_r=0.7\cdot T_c$ falls at a different absolute temperature for each method (because each has its own $T_c$) and the datasheet has no tabulated point at exactly that temperature, $P^{sat}(0.7\,T_c)$ is instead obtained by interpolating between the two nearest tabulated points (see Remark below). For NIST ($T_c=351.3\,\text{K}$, so $T_r=0.7\to-27.24\,^\circ\text{C}$), this falls between the $-30\,^\circ\text{C}$ ($P^{sat}=2.7344\,\text{bar}$) and $-10\,^\circ\text{C}$ ($P^{sat}=5.8263\,\text{bar}$) tabulated points. Although all three methods draw from this same Linde table, each computes a \emph{different} $\omega$, because each uses its own fitted $T_c$ and $P_c$ in the formula above:
\[
\omega_{\text{NIST}} = 0.2769, \qquad \omega_{\text{GCGP}} = 0.1711, \qquad \omega_{\text{SPGP}} = -0.2741.
\]
SPGP's negative value is not a physically meaningful acentric factor; it is a direct consequence of SPGP's fitted $T_c$ ($400.9\,\text{K}$) being far higher than R-32's actual critical temperature, which artificially depresses the computed reduced saturation pressure at $T_r=0.7$.
\begin{remark}
The interpolation of the Linde saturation-pressure data used to determine
$T_r=0.7$ saturation points is performed in $(1/T,\ \ln P)$ coordinates
rather than $(T,P)$ directly, following the Clausius-Clapeyron relation
\cite[sec.\ 7-2]{poling2001}, $d(\ln P)/d(1/T) \approx -\Delta H_{vap}/R$,
which is approximately linear over a limited temperature range and is the
same theoretical basis underlying standard vapor-pressure correlations
such as the Antoine equation.
\end{remark}
\begin{remark}
This Linde-anchored determination of $\omega$ presupposes the availability
of an experimental saturation-pressure curve for the fluid in question,
which holds for R-32 as a reference case but not, by construction, for the
novel and non-standard refrigerant candidates that GCGP and SPGP are
intended to screen. Neither the Ambrose-Walton correlation~\cite{ambrosewalton1989}
nor an equal-fugacity flash within the PR equation of state itself can
substitute for this step, since both require $\omega$ as an input --- through
the $\omega f^{(1)} + \omega^2 f^{(2)}$ terms in the former, and through $\kappa(\omega)$
in the latter --- rather than producing it. \textbf{Future novel candidates will explicitly require $P_{sat}$ at $T_r = \frac{T}{T_c} = 0.7$.}
\end{remark}
\subsubsection{Computation of Saturation Properties}\label{sec:ambrose_walton}
The Ambrose-Walton corresponding-states correlation~\cite{ambrosewalton1989}
is used to compute the saturation pressure $P^{sat}(T)$ once $\omega$ is
known:
\begin{subequations}
\begin{align}
    \ln P_r^{sat} &= f^{(0)}(T_r) + \omega f^{(1)}(T_r) + \omega^2 f^{(2)}(T_r) \\
    f^{(0)}(T_r) &= \frac{-5.97616\,\tau + 1.29874\,\tau^{1.5} - 0.60394\,\tau^{2.5} - 1.06841\,\tau^{5}}{T_r} \\
    f^{(1)}(T_r) &= \frac{-5.03365\,\tau + 1.11505\,\tau^{1.5} - 5.41217\,\tau^{2.5} - 7.46628\,\tau^{5}}{T_r} \\
    f^{(2)}(T_r) &= \frac{-0.64771\,\tau + 2.41539\,\tau^{1.5} - 4.26979\,\tau^{2.5} + 3.25259\,\tau^{5}}{T_r}
\end{align}
\end{subequations}
where $T_r = \frac{T}{T_c}$, $\tau = 1 - T_r$, and $P_r^{sat} = \frac{P^{sat}}{P_c}$.
\subsection{Caloric Properties for the Vanilla Python}\label{sec:vanilla_caloric}
Caloric properties predicted as part of the equation of state package include the Enthalpy $h$ and the entropy $s$. Both expression are computed as a sum of their ideal-gas and departure terms.
\begin{subequations}
    \begin{align}
    h(T,P) &= h^{ig}(T) + \Delta h^{dep}(T,P)\label{eq:vanilla_enthalpy}\\
    s(T,P) &= s^{ig}(T,P) + \Delta s^{dep}(T,P)\label{eq:vanilla_entropy}
    \end{align}
\end{subequations}

The ideal-gas terms are computed from the ideal-gas heat capacity, which in turn is computed from the Shomate Expression~\cite{shomate1954}:
\begin{equation}\label{eq:shomate_cp}
    C_p^{ig}(T) = A + Bt + Ct^2 + Dt^3 + \frac{E}{t^2}, \qquad t = \frac{T}{1000}
\end{equation}
The Eq.~\eqref{eq:shomate_cp} gives the expressions for the ideal-gas enthalpy and the entropy with respect to the IIR reference state convention: $T_{ref} =\SI{273.15}{\kelvin}$ and $P_{ref} = 1~\text{bar}$~\cite{ashrae2009}.
\begin{subequations}
\begin{align}
    H^{ig}(T) - H^{ig}(T_{ref}) &= \left[At + \frac{B}{2}t^2 + \frac{C}{3}t^3 + \frac{D}{4}t^4 - \frac{E}{t}\right]_{T_{ref}}^{T} \label{eq:ideal_enthalpy} \\
    S^{ig}(T,P) - S^{ig}(T_{ref},P_{ref}) &= \left[A\ln t + Bt + \frac{C}{2}t^2 + \frac{D}{3}t^3 - \frac{E}{2t^2}\right]_{T_{ref}}^{T} - R\ln\left(\frac{P}{P_{ref}}\right) \label{eq:ideal_entropy}
\end{align}
\end{subequations}
The departure functions for enthalpy and entropy follow the general
cubic-EoS expressions of Poling et al., Table 6-3~\cite{poling2001}:
\begin{subequations}
\begin{align}
    \frac{H^{ig}-H}{RT} &= -\frac{\Theta - T\dfrac{d\Theta}{dT}}{RT\sqrt{\delta^2-4\varepsilon}}\ln\left[\frac{2V+\delta-\sqrt{\delta^2-4\varepsilon}}{2V+\delta+\sqrt{\delta^2-4\varepsilon}}\right] + 1 - Z \\
    \frac{S^{ig}-S}{R} &= \frac{1}{R\sqrt{\delta^2-4\varepsilon}}\frac{d\Theta}{dT}\ln\left[\frac{2V+\delta-\sqrt{\delta^2-4\varepsilon}}{2V+\delta+\sqrt{\delta^2-4\varepsilon}}\right] - \ln\left[Z\left(1-\frac{b}{V}\right)\right]
\end{align}
\end{subequations}
For the PR equation of state specifically, $\Theta=a(T)$, $\dfrac{d\Theta}{dT}=\dfrac{da}{dT}$, $\delta=ub$, $\varepsilon=wb^2$, so $\sqrt{\delta^2-4\varepsilon}=b\sqrt{u^2-4w}\equiv b\sigma$. Substituting these and converting $V\to Z$ (via $V=\dfrac{ZRT}{P}$, $b=\dfrac{BRT}{P}$) gives:
\begin{subequations}
\begin{align}
    \Delta h^{dep}(T,P) &= H - H^{ig} = RT(Z-1) + \frac{T\dfrac{da}{dT} - a}{b\sigma}\ln\left[\frac{2Z+B(u+\sigma)}{2Z+B(u-\sigma)}\right]\label{eq:dh_dep} \\
    \Delta s^{dep}(T,P) &= S - S^{ig} = R\ln(Z-B) + \frac{1}{b\sigma}\frac{da}{dT}\ln\left[\frac{2Z+B(u+\sigma)}{2Z+B(u-\sigma)}\right]\label{eq:ds_dep}
\end{align}
\end{subequations}
where $\sigma \equiv \sqrt{u^2-4w}$. Since $u=2$ and $w=-1$ are universal
PR constants, $\sigma=\sqrt{8}=2\sqrt{2}$ is itself a universal constant,
the same for every fluid.
\section{Translation of the Vanilla Property Package to IDAES Custom Property Package}\label{sec:idaes_translation}
The custom property package re-declares the PR equation of state and Shomate ideal-gas correlation from Section~\ref{sec:pr_pvt} inside IDAES's generic property framework~\cite{idaes_cubic_eos_2026}. The underlying equations are not re-derived: the PR parameterization ($a(T)$, $b$, $\alpha(T)$, $\kappa$, and the universal constants $u=2$, $w=-1$, $\Omega_A=0.45724$, $\Omega_B=0.0778$, Eq.~\eqref{eq:PR} and the $\kappa$ relation, Section~\ref{sec:pr_pvt}) and the Shomate ideal-gas correlation ($C_p^{ig}$, $H^{ig}$, $S^{ig}$, Eqs.~\eqref{eq:shomate_cp}, \eqref{eq:ideal_enthalpy}, \eqref{eq:ideal_entropy}) are identical to the vanilla implementation. This equivalence was validated numerically, with the package's computed compressibility factor $Z$ matching the vanilla calculation to within $10^{-3}$.

The phase equilibrium computations differ in that they do not use the
Ambrose--Walton method. The property
package instead locates the phase boundary directly from the PR equation of state, via the equal-fugacity criterion for vapor-liquid
equilibrium~\cite[sec.\ 8-2]{poling2001}:
\begin{equation}\label{eq:equal_fugacity}
    f^{L}(T,P) = f^{V}(T,P)
\end{equation}
for a pure fluid. The fugacity coefficient is defined as $\phi\equiv \frac{f}{P}$. For a pure fluid it follows from the general cubic-EoS fugacity departure function of Poling et al., Table 6-3~\cite{poling2001}:
\begin{equation}
    \ln\left(\frac{f}{P}\right) = \frac{\Theta}{RT\sqrt{\delta^2-4\varepsilon}}\ln\left[\frac{2V+\delta-\sqrt{\delta^2-4\varepsilon}}{2V+\delta+\sqrt{\delta^2-4\varepsilon}}\right] - \ln\left[Z\left(1-\frac{b}{V}\right)\right] - (1-Z)
\end{equation}
Applying the same PR substitution used in Section~\ref{sec:vanilla_caloric} ($\Theta=a$,
$\delta=ub$, $\varepsilon=wb^2$, $\sqrt{\delta^2-4\varepsilon}=b\sigma$)
and converting $V\to Z$ gives:
\begin{equation}\label{eq:log_fug_coeff}
    \ln\phi \equiv \ln\left(\frac{f}{P}\right) = (Z-1) - \ln(Z-B) - \frac{A}{B\sigma}\ln\left[\frac{2Z+B(u+\sigma)}{2Z+B(u-\sigma)}\right]
\end{equation}
Substituting $f=P\phi$ into
Eq.~\eqref{eq:equal_fugacity}, the pressure cancels (both phases are
evaluated at the same $T,P$), so the condition reduces to equality of
the fugacity coefficients:
\begin{equation}\label{eq:equal_phi}
    \phi^{L}(T,P) = \phi^{V}(T,P)
\end{equation}
Writing each side explicitly via Eq.~\eqref{eq:log_fug_coeff}, evaluated
at the liquid root $Z_L$ and vapor root $Z_V$ of Eq.~\eqref{eq:cubic_eos}
respectively, Eq.~\eqref{eq:equal_phi} becomes:
\begin{equation}\label{eq:equal_fugacity_explicit}
\begin{aligned}
    (Z_L-1) - \ln(Z_L-B) - \frac{A}{B\sigma}\ln\left[\frac{2Z_L+B(u+\sigma)}{2Z_L+B(u-\sigma)}\right] \\
    = (Z_V-1) - \ln(Z_V-B) - \frac{A}{B\sigma}\ln\left[\frac{2Z_V+B(u+\sigma)}{2Z_V+B(u-\sigma)}\right]
\end{aligned}
\end{equation}
This is a single explicit equation in $Z_L$, $Z_V$, and hence in $T$ and
$P$ (through $A(T,P)$, $B(T,P)$), coupled with the requirement that
$Z_L$ and $Z_V$ are themselves both real roots of the same cubic,
Eq.~\eqref{eq:cubic_eos}, at that $T,P$. For a pure fluid, fixing $T$
and solving this coupled system for the $P$ that satisfies
Eq.~\eqref{eq:equal_fugacity_explicit} gives $P^{sat}(T)$ directly from
the equation of state itself --- no external correlation is used, in
contrast to the vanilla implementation's reliance on Ambrose-Walton.

IDAES's cubic equation-of-state implementation computes the departure
enthalpy and entropy via its own \texttt{enth\_mol\_phase} and
\texttt{entr\_mol\_phase} functions. Each independently implements the same
general cubic-EoS departure formulas of Poling et al., Table
6-3~\cite{poling2001} used in Section~\ref{sec:vanilla_caloric} --
confirmed directly from the source code, whose comments cite the same
underlying reference (``pg.\ 120'' and ``pg.\ 102, or Sandler, 4th ed.,
pg.\ 208,'' respectively) -- so they reduce to exactly the same
expressions already derived there, Eqs.~\eqref{eq:dh_dep}--\eqref{eq:ds_dep}:
\begin{subequations}
\begin{align}
    \Delta h^{dep}(T,P) &= RT(Z-1) + \frac{T\dfrac{da}{dT} - a}{b\sigma}\ln\left[\frac{2Z+B(u+\sigma)}{2Z+B(u-\sigma)}\right] \\
    \Delta s^{dep}(T,P) &= R\ln(Z-B) + \frac{1}{b\sigma}\frac{da}{dT}\ln\left[\frac{2Z+B(u+\sigma)}{2Z+B(u-\sigma)}\right]
\end{align}
\end{subequations}
These are computed independently of the fugacity coefficient,
Eq.~\eqref{eq:log_fug_coeff} -- IDAES does not derive one from the other
internally. Rather, the identity $\ln\phi = G^{dep}/(RT) =
(\Delta h^{dep}-T\Delta s^{dep})/(RT)$ (from $G=h-TS$~\cite{smithvanness2018})
is a thermodynamic consistency relation that all three internally-computed
quantities must satisfy, and do, since each is a standard Table 6-3 form
evaluated at the same $Z$, $A$, $B$.

No separate departure-function or fugacity-coefficient option is exposed in
the property-package configuration; selecting \texttt{equation\_of\_state: Cubic}
with \texttt{CubicType.PR} is sufficient for IDAES to build and evaluate all of the
above internally, using the same $(u,w,\Omega_A,\Omega_B)$ and $(T_c,P_c,\omega)$
already fixed in Section~\ref{sec:pr_pvt}.

\subsection{Verifying the IDAES Property Package Against the Vanilla Implementation}\label{sec:idaes_verification}
The equivalence between the vanilla and IDAES implementations was validated
numerically: for each method (NIST, GCGP, SPGP), the compressibility factor
$Z$ computed by the custom IDAES property package matched the vanilla
calculation to within a tolerance of $10^{-3}$. As an independent
cross-check on the acentric factor specifically, IDAES's
\texttt{CoolPropWrapper} database lookup was queried directly for R-32 and
returned $\omega=0.2769$, matching the Linde-interpolated value used in
Section~\ref{sec:acentric_factor} to four decimal places.
\section{Single-Stage Vapor-Compression Cycle Thermodynamic Balances}\label{sec:cycle_balances}
DVCT models a single-stage vapor-compression cycle as a closed loop of four
units -- evaporator, compressor, condenser, expansion valve -- carrying a
single pure-fluid molar flow rate $\dot n$, with no splitting or mixing
anywhere in the loop. Each unit is modeled as a steady-state, single-inlet
single-outlet control volume, for which the general first-law energy
balance (negligible kinetic and potential energy) is:
\begin{equation}\label{eq:energy_balance}
    \dot n\, h_{out} = \dot n\, h_{in} + \dot Q + \dot W
\end{equation}
where $\dot Q$ is heat added to the fluid and $\dot W$ is shaft work done
on the fluid.

\paragraph{Evaporator.} No shaft work ($\dot W=0$); Eq.~\eqref{eq:energy_balance}
reduces to the refrigeration effect:
\begin{equation}\label{eq:q_evap}
    \dot Q_{evap} = \dot n\left(h_{evap,out} - h_{evap,in}\right)
\end{equation}
The evaporator's operating point is set directly by a specified saturation
temperature setpoint $T_{evap,set}$ -- the target refrigeration
temperature -- imposed as an equality constraint on the evaporator
outlet's saturation (bubble-point) temperature:
\begin{equation}\label{eq:evap_setpoint}
    T_{sat}(P_{evap}) = T_{evap,set}
\end{equation}
Since $T_{sat}(P)$ is the same monotonic saturation curve established by
the equal-fugacity condition in Section~\ref{sec:idaes_translation}, fixing
$T_{evap,set}$ is what actually determines $P_{evap}$ -- the low-side
pressure is a consequence of the temperature setpoint, not specified
independently. Physically, $T_{evap,set}$ sits below the temperature of
the space being cooled by an evaporator approach temperature, the same
convention used on the condenser side (below). Standard frozen-food
storage rooms are held at $-18$ to $-20^\circ$C; with an 8-10$^\circ$C
evaporator approach, this places $T_{evap,set}=-29^\circ$C squarely
within the range needed for freezer duty, not an unusually low choice.
The outlet is additionally, optionally constrained to a specified
superheat $\Delta T_{sh}$ above this same saturation temperature:
\begin{equation}\label{eq:superheat}
    T_{evap,out} \ge T_{sat}(P_{evap}) + \Delta T_{sh}, \qquad \Delta T_{sh}=0 \text{ by default}
\end{equation}

\paragraph{Compressor.} Adiabatic ($\dot Q=0$); Eq.~\eqref{eq:energy_balance}
gives the actual compression work:
\begin{equation}\label{eq:w_comp}
    \dot W_{comp} = \dot n\left(h_{comp,out} - h_{comp,in}\right)
\end{equation}
The isentropic efficiency $\eta_s$ relates this to the work required for a
hypothetical reversible (isentropic) compression to the same outlet
pressure, $h_{isen}$ evaluated at $s_{isen}=s_{comp,in}$, $P_{isen}=P_{comp,out}$:
\begin{equation}\label{eq:eta_isen}
    \eta_s \equiv \frac{h_{isen} - h_{comp,in}}{h_{comp,out} - h_{comp,in}}
\end{equation}

\paragraph{Condenser.} No shaft work; Eq.~\eqref{eq:energy_balance} gives the
heat rejected:
\begin{equation}\label{eq:q_cond}
    \dot Q_{cond} = \dot n\left(h_{cond,in} - h_{cond,out}\right)
\end{equation}
Unlike the evaporator, the condenser's operating point is not set by a
direct temperature setpoint. Instead it is tied to the ambient
temperature $T_{amb}$ -- an external boundary condition representing the
surroundings the condenser rejects heat into, not a state internal to the
refrigerant loop -- via a specified approach temperature $\Delta T_{app}$,
the minimum temperature difference required to drive heat rejection into
that ambient:
\begin{equation}\label{eq:subcool_approach}
    T_{sat}(P_{cond}) = T_{amb} + \Delta T_{app}
\end{equation}
This is the same approach-temperature convention noted for the evaporator
above, applied on the rejection side instead of the absorption side. As
with the evaporator, this fixes $P_{cond}$ through the same saturation
curve. The outlet is additionally, optionally constrained to a specified
subcooling $\Delta T_{sc}$ below this saturation temperature:
\begin{equation}\label{eq:subcool}
    T_{cond,out} \le T_{sat}(P_{cond}) - \Delta T_{sc}, \qquad \Delta T_{sc}=0 \text{ by default}
\end{equation}

\paragraph{Expansion valve.} Adiabatic and no shaft work ($\dot Q=\dot W=0$);
Eq.~\eqref{eq:energy_balance} reduces to the isenthalpic throttling condition:
\begin{equation}\label{eq:valve_isenthalpic}
    h_{valve,out} = h_{valve,in}
\end{equation}

\paragraph{Mass balance.} Since the loop is a single pure-fluid stream with
no splitting or mixing, $\dot n$ is the same constant at every point in the
cycle.

\paragraph{Coefficient of performance.} The cycle's COP is defined as the
ratio of refrigeration effect to compressor work, Eqs.~\eqref{eq:q_evap}
and~\eqref{eq:w_comp}; since $\dot n$ is common to both and cancels:
\begin{equation}\label{eq:cop}
    COP \equiv \frac{\dot Q_{evap}}{\dot W_{comp}} = \frac{h_{evap,out}-h_{evap,in}}{h_{comp,out}-h_{comp,in}}
\end{equation}

\paragraph{Carnot efficiency.} The second law bounds the COP of any cycle
operating between the evaporator and condenser saturation temperatures
(absolute temperature, $T_{evap,set}$ and $T_{sat}(P_{cond})$ from
Eqs.~\eqref{eq:evap_setpoint} and~\eqref{eq:subcool_approach}):
\begin{equation}\label{eq:cop_carnot}
    COP_{Carnot} = \frac{T_{evap,set}}{T_{sat}(P_{cond}) - T_{evap,set}}
\end{equation}
and the cycle's second-law efficiency is the ratio of the actual COP,
Eq.~\eqref{eq:cop}, to this bound:
\begin{equation}\label{eq:eta_second_law}
    \eta_{II} \equiv \frac{COP}{COP_{Carnot}}
\end{equation}

\paragraph{Operating conditions used in this work.} The cycle's ideal-cycle
modeling approach -- isentropic compression ($\eta_s\approx1$), no superheat or subcooling -- follows Raskar and Mutalikdesai~\cite{raskarmutalikdesai2016}. The evaporator setpoint, $T_{evap,set}=-29^\circ$C, targets standard freezer duty (a $-18$ to
$-20^\circ$C storage room with an 8-10$^\circ$C evaporator approach); the ambient sweep is restricted to $T_{amb}\in[10,25]^\circ$C, since this deep-evaporator duty becomes solver-limited at higher ambient (a limitation of this operating regime, not of R-32 or the property models themselves). The condenser approach is $\Delta T_{app}=9^\circ$C, and the compressor pressure ratio is bounded at $r_p\le10$.

\begin{remark}
Physically, $T_{evap,set}$ sits below the temperature of the space being cooled by an evaporator approach temperature, the same convention used for the condenser's approach to ambient, Eq.~\eqref{eq:subcool_approach}.
Standard frozen-food storage rooms are held at $-18$ to $-20^\circ$C. Applying the same 8-10$^\circ$C approach used elsewhere in this cycle reproduces the chosen setpoint from either direction: forward, $T_{box} -
\Delta T_{approach}$ gives $-26$ to $-30^\circ$C; backward, $T_{evap,set} +
\Delta T_{approach} = -29 + 8$ to $-29 + 10 = -21$ to $-19^\circ$C, landing almost exactly on the standard box-temperature range. So
$T_{evap,set}=-29^\circ$C corresponds to ordinary freezer duty, not an unusually low or arbitrary choice.
\end{remark}

\section{Algebraic (Sequential) Evaluation of the Ideal Cycle}\label{sec:algebraic_cycle}

Sections~\ref{sec:cycle_balances}'s four balances -- Eqs.~\eqref{eq:q_evap},
\eqref{eq:w_comp}, \eqref{eq:q_cond}, \eqref{eq:valve_isenthalpic} -- together
with the equal-fugacity saturation condition of
Section~\ref{sec:idaes_translation} and the caloric relations of
Section~\ref{sec:vanilla_caloric}, are the same physics IDAES poses as one
large simultaneous nonlinear system and solves with IPOPT across all four
unit models at once. This section shows that, under the operating
conditions actually used in this work -- $T_{evap,set}$ and $T_{amb}$
specified, $\Delta T_{sh}=\Delta T_{sc}=0$ (the defaults in
Eqs.~\eqref{eq:superheat}, \eqref{eq:subcool}) -- the identical system of
equations can instead be evaluated as a single deterministic forward pass,
with no simultaneous solve and only one one-dimensional root find. This
matters specifically for a parameter sensitivity study over $T_c$ and
$P_c$: a simultaneous IPOPT solve carries path-dependent convergence
behavior, a nonzero tolerance floor, and warmstart-sensitive failures that
are unrelated to the underlying physics and can be mistaken for genuine
sensitivity. The sequential evaluation below removes that ambiguity by
construction.

\subsection{Why the System Is Sequential Rather Than Simultaneous}\label{sec:why_sequential}

With $T_{evap,set}$ and $T_{amb}$ fixed, Eqs.~\eqref{eq:evap_setpoint} and
\eqref{eq:subcool_approach} already fix $P_{evap}$ and $P_{cond}$
independently of every other unknown in the cycle, since $T_{sat}(P)$ is a
function of $P$ alone (Sections~\ref{sec:ambrose_walton},
\ref{sec:idaes_translation}). With $\Delta T_{sh}=\Delta T_{sc}=0$, the
evaporator outlet and condenser outlet are pinned exactly to this
saturation curve, so their states ($h$, $s$) are likewise fully determined
without reference to any other unit. The only unit whose outlet state is
not immediately fixed by a specified temperature is the compressor, whose
outlet is instead constrained implicitly through $\eta_s$
(Eq.~\eqref{eq:eta_isen}) and the isentropic condition; and the valve
outlet, Eq.~\eqref{eq:valve_isenthalpic}, is fixed once the condenser
outlet enthalpy is known. Because each unit's governing equation involves
at most one quantity not already determined by the unit immediately
upstream of it, the four balances can be evaluated once, in order, rather
than posed as one system in all unknowns simultaneously.

\subsection{State-by-State Evaluation}\label{sec:state_by_state}

\paragraph{Evaporator outlet.} From Eq.~\eqref{eq:evap_setpoint},
\begin{equation}\label{eq:seq_evap_out}
    P_{evap} = P^{sat}(T_{evap,set}), \qquad
    h_{evap,out} = h(T_{evap,set}, P_{evap}), \qquad
    s_{evap,out} = s(T_{evap,set}, P_{evap})
\end{equation}
evaluated at the vapor root $Z_V$ via Eqs.~\eqref{eq:vanilla_enthalpy},
\eqref{eq:vanilla_entropy}. This is a direct evaluation: no unknown remains
to be found.

\paragraph{Compressor outlet.} The isentropic outlet temperature $T_{isen}$
is the unique root of
\begin{equation}\label{eq:seq_isentropic_root}
    s\big(T_{isen}, P_{cond}\big) - s_{evap,out} = 0
\end{equation}
where $P_{cond}=P^{sat}(T_{sat}(P_{cond}))$ follows from
Eq.~\eqref{eq:subcool_approach} below and $s(T,P_{cond})$ is evaluated at
the vapor root $Z_V(T,P_{cond})$. Because $s(T,P_{cond})$ increases
monotonically with $T$ along the vapor branch, Eq.~\eqref{eq:seq_isentropic_root}
has exactly one root and can be found by a bracketed one-dimensional search
(e.g.\ bisection); this is the only implicit step in the entire sequential
evaluation, and involves a single scalar unknown rather than the coupled
system IDAES solves. With $h_{isen}=h(T_{isen},P_{cond})$, the actual
outlet enthalpy follows from Eq.~\eqref{eq:eta_isen} rearranged:
\begin{equation}\label{eq:seq_comp_out}
    h_{comp,out} = h_{evap,out} + \frac{h_{isen} - h_{evap,out}}{\eta_s}
\end{equation}

\paragraph{Condenser outlet.} From Eq.~\eqref{eq:subcool_approach},
\begin{equation}\label{eq:seq_cond_out}
    P_{cond} = P^{sat}\!\left(T_{amb}+\Delta T_{app}\right), \qquad
    h_{cond,out} = h\!\left(T_{amb}+\Delta T_{app},\, P_{cond}\right), \qquad
    s_{cond,out} = s\!\left(T_{amb}+\Delta T_{app},\, P_{cond}\right)
\end{equation}
now evaluated at the liquid root $Z_L$. As with the evaporator outlet,
this is a direct evaluation, independent of the compressor and valve
states.

\paragraph{Valve outlet.} From Eq.~\eqref{eq:valve_isenthalpic},
\begin{equation}\label{eq:seq_valve_out}
    h_{valve,out} = h_{cond,out}
\end{equation}
a direct assignment rather than a calculation. The resulting state lies
within the two-phase region at $P_{evap}$; its vapor quality $x$ follows
by flashing at known enthalpy between the saturated-liquid and
saturated-vapor enthalpies at $T_{evap,set}$ -- the latter already
available from Eq.~\eqref{eq:seq_evap_out}, the former obtained the same
way at the liquid root:
\begin{equation}\label{eq:seq_quality}
    x = \frac{h_{valve,out} - h_{f}(T_{evap,set})}{h_{g}(T_{evap,set}) - h_{f}(T_{evap,set})}, \qquad
    s_{valve,out} = s_f(T_{evap,set}) + x\big[s_g(T_{evap,set}) - s_f(T_{evap,set})\big]
\end{equation}
where $h_f,\,s_f$ and $h_g,\,s_g$ denote the saturated liquid ($Z_L$) and
saturated vapor ($Z_V$) properties at $T_{evap,set}$, $P_{evap}$. No
additional root-finding is required, since $h_{valve,out}$ is already
known from Eq.~\eqref{eq:seq_valve_out}.

\subsection{Coefficient of Performance}\label{sec:seq_cop}

Substituting $h_{evap,in}=h_{valve,out}$ and
$h_{comp,in}=h_{evap,out}$ into Eq.~\eqref{eq:cop} recovers the same COP
definition already established in Section~\ref{sec:cycle_balances}, now
built entirely from the direct and single-root-find evaluations above:
\begin{equation}\label{eq:seq_cop}
    COP = \frac{h_{evap,out} - h_{valve,out}}{h_{comp,out} - h_{evap,out}}
\end{equation}
identical in form to Eq.~\eqref{eq:cop}. Because COP is a ratio of
enthalpy \emph{differences}, the IIR reference-state constant fixed in
Eq.~\eqref{eq:ideal_enthalpy} cancels exactly in
Eq.~\eqref{eq:seq_cop}; no reference-state anchoring step is required for
the sequential evaluation, unlike quantities (e.g.\ absolute $h$ or $s$)
that are compared directly against a datasheet.

\subsection{Summary}\label{sec:seq_summary}

\begin{center}
\begin{tabular}{lll}
\textbf{Quantity} & \textbf{Type of evaluation} & \textbf{Governing equation} \\
\hline
$P_{evap}$, $P_{cond}$ & closed-form & Eqs.~\eqref{eq:seq_evap_out}, \eqref{eq:seq_cond_out} \\
$h_{evap,out}$, $h_{cond,out}$ & closed-form & Eqs.~\eqref{eq:vanilla_enthalpy}--\eqref{eq:vanilla_entropy} \\
$T_{isen}$ (compressor) & 1-D root find & Eq.~\eqref{eq:seq_isentropic_root}, monotone in $T$ \\
$h_{comp,out}$ & closed-form, given $T_{isen}$ & Eq.~\eqref{eq:seq_comp_out} \\
$h_{valve,out}$, $x$ & closed-form & Eqs.~\eqref{eq:seq_valve_out}, \eqref{eq:seq_quality} \\
\hline
Full IDAES flowsheet & simultaneous solve (IPOPT) & Section~\ref{sec:idaes_translation}, not used here \\
\end{tabular}
\end{center}

The sequential evaluation above uses exactly the same equation-of-state,
departure-function, and saturation relations already established in
Sections~\ref{sec:pr_pvt}--\ref{sec:idaes_translation}; nothing about the
underlying physics changes. What changes is only the numerical procedure
by which the four states are found: one deterministic forward pass with a
single bracketed root find, in place of one simultaneous nonlinear system
solved across all four unit models together.

\bibliographystyle{plain}
\bibliography{ref}
\end{document}
